\chapter*{ЗАКЛЮЧЕНИЕ}
\addcontentsline{toc}{chapter}{ЗАКЛЮЧЕНИЕ}

Являясь субъектом многообразных явлений и процессов, мы нуждаемся в разнообразных технологиях наблюдения,
анализа и прогнозирования явлений, описываемых большим объемом информации.
Одним из средств достижения данных практических целей является использование временных рядов.
Реализуя поставленные в работе задачи, мы:
\begin{itemize}
    \item проанализировали предметную область «временных рядов»;
    \item сформулировали критерии сравнения методов хранения и прогнозирования временных рядов;
    \item описали и классифицировали методы хранения и прогнозирования временных рядов.
\end{itemize}

В работе выделены параметры и произведена классификация систем хранения разных видов с целью определения их области применения для практического использования следующим образом:

\begin{itemize}
    \item Cassandra --- метод с автоматизированной системой сортировки данных по времени и удаления устаревшей информации;
    \item Cassandra и OpenTSDB --- системы, отвечающие требованиям отказоустойчивости, масштабируемости и приоритетности скорости записи данных по отношению к скорости их изменения;
    \item OpenTSDB --- система, в которой есть инструменты для интеллектуального анализа временных рядов и встроенный модуль для визуализации данных;
    \item Файловая система дает возможность максимально сжать информацию.
\end{itemize}

Методы прогнозирования временных рядов дают возможность предугадать динамику отслеживаемых процессов.
В работе были рассмотрены методы построения статистических моделей (экспоненциального сглаживания, регрессионные, авторегрессионные)
и структурных моделей (на цепях Маркова и на классификационно-регрессионных деревьях).

С точки зрения необходимости специальной обработки данных перед прогнозированием и возможности учета влияния внешних категориальных факторов можно разделить модели на группы:
\begin{itemize}
    \item статистические модели, не требующие специальной обработки данных и не учитывающие влияние внешних категориальных факторов;
    \item структурные модели, которые требуют специальной обработки данных и учитывают влияние внешних категориальных факторов.
\end{itemize}

Учитывая количество учитываемых внешних дискретных факторов, классифицируем методы следующим образом:
\begin{itemize}
    \item без учета внешних факторов: модель экспоненциального сглаживания;
    \item с учетом одного дискретного внешнего фактора: простая линейная и нелинейная регрессии и авторегрессионная модель по выборке максимального подобия;
    \item с учетом любого количества внешних факторов: структурные модели, модель множественной регрессии и авторегрессионные модели.
\end{itemize}

Кроме того, все рассмотренные модели имеют различную область применения.